\subsubsection{Vectorisation des colonnes sur processeur}
La boucle interne de NvPD calcule une ligne à partir de la table des dernières occurrences. L'instruction SIMD applique une même opération à plusieurs valeurs (\emph{single instruction, multiple data}). Le branchement sur \(k=0\), \(k=j\) ou \(0<k<j\) peut être remplacé par le minimum de l'algorithme~\ref{alg:nvpd}. Quand \(k=j\), le second terme vaut \(P_{j-1}\), le coût d'une correspondance. Le cas sans correspondance suit la borne établie pour l'algorithme~\ref{alg:nvpd}. Une instruction SIMD peut alors traiter plusieurs colonnes qui choisissent des prédécesseurs différents.
\chart{cpu/variantes_cpu.png}{Chaque case donne le gain moyen de trois blocs appariés. En dessous figurent l'écart-type et la demi-largeur de l'IC à 95\,\%. Le vert indique un gain supérieur à l'unité, le rose un ralentissement.}{fig:simd}{bench/cpu/openmp/profile_cpu_driver.cpp}
Les essais de la figure~\ref{fig:simd} utilisent GCC~11.5 avec \texttt{-O3 -march=native} sur EPYC~9654. La moyenne des rapports \(T_\mathrm{initial}/T_\mathrm{SIMD}\) porte sur trois répétitions appariées. Le compilateur vectorise la boucle sur 32 octets, soit huit entiers par instruction vectorielle. À un cœur, le gain vaut \(5{,}852\pm0{,}009\) sur la forme \(256\times10^6\). Le panneau b) de la figure~\ref{fig:million} compare ensuite la boucle initiale et la boucle SIMD dans cinq blocs appariés, sur les mêmes huit cœurs et la même entrée par taux. À 2\,\%, le temps passe de \(321{,}572\pm0{,}385\) à \(51{,}655\pm0{,}368\)\,s. Sur \(10^6\times1\,000\,033\) cases, le débit moyen du noyau SIMD atteint \(19{,}36\pm0{,}14\)\,GCUPS. Pour 2, 5 et 10\,\% d'éditions injectées, les gains appariés valent \(6{,}226\pm0{,}043\), \(6{,}192\pm0{,}038\) et \(6{,}206\pm0{,}049\).
\chart{cpu/cout_instructions.png}{Deux processus par configuration mesurent les instructions retirées par case et les erreurs de prédiction pour mille cases. Sous chaque moyenne figurent l'écart-type et la demi-largeur de l'IC à 95\,\%.}{fig:instructions}{bench/cpu/openmp/profile_cpu.py}
La commande \texttt{perf stat} de Linux (version~6.9.10) relève les compteurs matériels de l'EPYC~\cite{perfstat}. La figure~\ref{fig:instructions} rapporte ces compteurs au travail dense \(W=mn\). Sur un seul cœur et \(20\,000^2\) cases, la reformulation ramène les instructions retirées de 21,76 à 2,13 par case. Les cycles comprennent aussi les dépendances de données et les attentes mémoire.

\synthese{Sur le carré de 20\,000 bases, le gain de la boucle vectorisée atteint \(2{,}62\pm0{,}53\) sur un cœur. La baisse des instructions retirées accompagne ce gain, même si les lectures indirectes restent présentes.}

\subsubsection{Empreinte de l’index et localité sur processeur}
La taille de l'index détermine la part du jeu de données qui tient en cache. Pour \(q\) lettres distinctes et \(n\) colonnes, l'index et les deux lignes du noyau de coût seul occupent
\[
 M_{\mathrm{NvPD}}=4(q+2)(n+1)\quad [\mathrm{bytes}].
\]
Les deux lignes sont contiguës. Le dernier prédécesseur varie avec la lettre. Le prototype compact code l'écart dans un octet lorsque sa valeur et les sentinelles tiennent dans cette représentation. La conversion garde provisoirement l'index entier. Sur ces processeurs, huit cœurs partagent 32\,Mio de cache L3. Chaque cœur dispose de 32\,Kio de L1 de données et de 1\,Mio de L2.
\chart{cpu/memoire_cpu.png}{À huit cœurs, les barres V1 (gauche) et V2 (droite) séparent les scores de l'index pour chaque matrice. Le trait donne le pic des tableaux et le losange le plus grand RSS observé. Le panneau b) rapporte les défauts L1D à mille cases.}{fig:mem}{bench/cpu/openmp/profile_cpu.py}
Sur la forme \(64\times4\,194\,304\), la figure~\ref{fig:mem} distingue le maximum d'espace de travail, qui passe de 96 à 80\,Mio, du RSS maximal observé parmi les processus. Sur ce grand rectangle à huit cœurs, les défauts L1 de données passent de 200,12 à 158,81 pour mille cases entre la version SIMD et sa variante compacte (deux processus de comptage par point). Un défaut L1 peut être servi par le L2 ou le L3.

\synthese{Sur la forme la plus large, l'index compact ramène le maximum des RSS observés de 101,6 à 85,5\,Mio. L'économie de mémoire dépend de la représentation et ne suffit pas à prédire la vitesse d'une ligne.}

\subsubsection{Décodage des fichiers chromosomiques}\label{sec:lecture}
Lorsque deux versions d'un chromosome diffèrent peu, la durée de lecture peut dépasser celle du calcul des fronts. Le banc isole la conversion du FASTA en deux chaînes de bases. Pour \(B\) octets de fichier et un temps de lecture \(T\), le débit vaut \(B/T\). Le processus relit le même fichier après échauffement. Cette mesure comprend le décodage et les allocations, ainsi que les copies depuis le cache de fichiers de Linux.

Le premier lecteur valide et convertit chaque lettre. Le deuxième ajoute d'un bloc les lignes composées uniquement de lettres A, C, G et T majuscules. Les autres lignes passent par une table de conversion qui conserve les symboles ambigus et ignore les espaces. Le troisième lecteur teste seize caractères à la fois avant de copier la ligne. La figure~\ref{fig:fasta} compare des variantes qui rendent les mêmes caractères sur les six chromosomes.
\chart{cpu/lecture_fasta_chromosomes.png}{Chaque abscisse représente un chromosome. Les compteurs de lecture sur EPYC (a à c) précèdent les rapports de temps appariés et les mesures d'allocation et de mémoire (d à f).}{fig:fasta}{bench/human/chromosome-panel/blocks/index.py}

Pour le chromosome~7, le traitement par ligne ramène les instructions retirées de 25,7 à 5,05 milliards. Les défauts de lecture L1 passent de 16,8 à 16,0 millions. Le lecteur exécute ainsi moins d'instructions, tandis que le nombre absolu de défauts varie peu. Sur le chromosome~1, le traitement explicite par blocs ramène ce compteur de 7,78 à 2,93 milliards sans modifier les 54 demandes d'allocation du processus pendant la lecture.
Le dernier essai réserve la capacité de chaque chaîne dès la lecture de son en-tête. L'index FAI\footnote{Le format \href{https://www.htslib.org/doc/faidx.html}{faidx de HTSlib} donne la longueur, la position et la largeur des lignes de chaque séquence.} fournit cette longueur. Dans neuf blocs distincts, le nombre d'appels à \texttt{new} passe de 53 à 15 pour le chromosome~1. Les octets cumulés demandés passent de 1\,006,64 à 496,79\,Mo. Sur le chromosome~20, le pic RSS moyen passe de 185,89 à 129,00\,Mio. La réservation évite de copier les caractères déjà stockés lorsque la chaîne grandit.

\synthese{Sur le chromosome~1, le lecteur par blocs passe de \(0{,}945\pm0{,}006\) à \(0{,}657\pm0{,}005\)\,s avec l'index. Le gain apparié vaut 1,44 (IC à 95\,\% de 1,43 à 1,45).}

\subsection{Traduction des adresses sur chromosomes entiers}\label{sec:pages-genomiques}
Le solveur relit les entrées à chaque alignement. Une entrée du cache de traduction couvre 512 fois plus de bases avec une page de 2\,Mio qu'avec une page de 4\,Kio, à raison d'un octet par base. Le banc place les deux séquences dans une zone anonyme alignée et mesure séparément le temps de copie \(A_p\) et le temps du score \(t_p\) pour chaque taille de page \(p\). Sous l'hypothèse d'un coût constant par score, le temps de \(k\) requêtes sur cette zone vaut
\[
                T_p(k)=A_p+k t_p .
\]
Sur un cœur EPYC, chaque chromosome passe par neuf blocs en ordre aléatoire, avec un appel d'échauffement par configuration. Linux confirme que les grandes pages couvrent la totalité de la zone traitée.\footnote{Le champ \texttt{AnonHugePages} de \texttt{/proc/self/smaps} mesure cette couverture. Voir la \href{https://docs.kernel.org/admin-guide/mm/transhuge.html}{documentation du noyau Linux}.} Edlib retrouve les distances de 176, 5 et 30 éditions sur les chromosomes 1, 18 et 20. La figure~\ref{fig:pages-genomiques} confronte les compteurs de traduction et le temps cumulé du calcul.
\chart[.81\linewidth]{cpu/pages_genomiques.pdf}{Les symboles donnent les moyennes des blocs, les traits épais un écart-type et les traits fins l'intervalle de Student à 95\,\%. La bande du modèle représente cet intervalle sur neuf blocs. Les séries chronométrées utilisent cinq blocs distincts et incluent la mise en mémoire des entrées.}{fig:pages-genomiques}{bench/cpu/huge-pages/trial.cpp}

Sur le chromosome~1, les grandes pages ramènent les demandes de parcours des tables de \(121\,594\pm6\) à \(53{,}5\pm5{,}9\) par score pour le solveur par fragments. Les défauts de lecture en cache L1 passent de 7,819 à 7,830 millions. Le débit de remplissage du cœur depuis la mémoire passe de \(6{,}27\pm0{,}01\) à \(8{,}04\pm0{,}03\)\,Go/s.\footnote{Le débit vaut \(64F/T\) pour \(F\) lignes reçues en \(T\) secondes. Le compteur Zen~4 inclut les demandes du cœur et le préchargement. Les écritures et les autres cœurs échappent à cette mesure. Voir la \reposource{results/profile/huge-pages/839316/events.txt}.} Un score prend alors \(18{,}80\pm0{,}04\)\,ms contre \(18{,}14\pm0{,}05\)\,ms avec les pages de 4\,Kio. Pour WFA2, le temps correspondant passe de \(89{,}28\pm0{,}15\) à \(87{,}47\pm0{,}08\)\,ms.

La copie des 476\,Mio du mappage prend \(211{,}85\pm1{,}44\)\,ms en petites pages et \(64{,}21\pm0{,}15\)\,ms en grandes pages. L'intersection des deux droites estimées se situe à 223 requêtes. À 192 requêtes effectivement chronométrées, l'écart de temps entre les deux tailles de page a un intervalle à 95\,\% de \(-83\) à \(+49\)\,ms. Le seuil dépend du coût du score sur les entrées déjà en mémoire.

\synthese{Avec la copie, les pages de 2\,Mio économisent \(147\pm2\)\,ms pour un alignement du chromosome~1. Sur 384 alignements avec une seule copie, les pages de 4\,Kio économisent \(146\pm36\)\,ms.}


\subsubsection{Passage à l’échelle en mémoire partagée}
Le passage à l'échelle forte répartit un travail constant entre \(p\) cœurs. Le passage à l'échelle faible garde le nombre de cases par cœur constant~\cite{amdahl,gustafson}. Pour une matrice de travail \(W\), les indicateurs valent
\[
 S_p=\frac{T_1(W)}{T_p(W)},\qquad E_p=\frac{S_p}{p},\qquad
 E_p^{\mathrm{weak}}=\frac{T_1(W_1)}{T_p(pW_1)}.
\]
La série forte utilise \(20\,000^2\) cases. La série faible conserve 4\,096 lignes pour \(4\,096p\) colonnes. Elle fixe donc le travail à 16\,777\,216 cases par cœur et élargit les tableaux avec \(p\).
\chart{cpu/echelle_cpu.png}{Les courbes donnent l'accélération forte et l'efficacité faible de NvPD selon le nombre de cœurs physiques de l'EPYC.}{fig:echellecpu}{bench/cpu/openmp/profile_cpu.sbatch}
À taille fixe, huit cœurs prennent \(27{,}96\pm0{,}09\)\,ms contre \(31{,}28\pm0{,}17\)\,ms à seize sur la figure~\ref{fig:echellecpu}. Au dernier placement, les cœurs occupent deux groupes de cache et deux domaines NUMA.\footnote{La distance entre les cœurs et les contrôleurs mémoire dépend du domaine NUMA auquel chacun appartient.} Les synchronisations entre lignes et les accès mémoire interviennent dans le même temps mesuré.

\synthese{Sur le carré de 20\,000 bases de côté, huit cœurs atteignent 73,8\,\% d’efficacité forte. À seize cœurs, la répartition sur deux groupes de cache s’accompagne d’une hausse du temps moyen.}
